\begin{titlepage}
\thispagestyle{empty}
\centering

\begin{center}
\includegraphics[width=0.18\linewidth]{uni\_logo.png}
\end{center}
\begin{otherlanguage}{arabic}
{\small جامعة العاصمة (جـامـعـة حـلـوان سابقاً)\par}
{\small كـليـة الحـاسـبـات والـذكـاء الاصطناعي\par}
{\small قـسـم الذكاء الاصطناعي\par}

\vspace{8mm}
{\Large\bfseries تايم لنز:\\
تطبيق للتعرُّف علي القطع الأثرية\\
وإجابة الأسئلة مُعزَّز بالاسترجاع\\
 للمتحف المصري الكبير
\par}
\vspace{8mm}

رسالة مشروع تخرج مقدمة من:\\
\vfill
\begin{tabular}{@{}ll@{}}
	{\large روان هشام \par} & {\large \textenglish{[20220167]}\par}\\
	{\large علي أشرف  \par} & {\large \textenglish{[20220296]}\par}\\
	{\large عمر أحمد  \par} & {\large \textenglish{[20220310]}\par}\\
	{\large عمر وجيه  \par} & {\large \textenglish{[20220324]}\par}\\
	{\large عمرو أحمد \par} & {\large \textenglish{[20210632]}\par}\\
	{\large ملك علاء   \par} & {\large \textenglish{[20220503]}\par}\\
\end{tabular}
\vfill
رسالة مقدمة ضمن متطلبات الحصول على درجة البكالوريوس في الحاسبات والذكاء الاصطناعي، بقسم الذكاء الاصطناعي، كلية الحاسبات والذكاء الاصطناعي، جامعة العاصمة (جامعة حلوان سابقاً)\\
\vfill
تحت إشراف:\\
{\large د. إسلام جمال\par}
\vfill
يونيو \textenglish{2026}
\end{otherlanguage}
\end{titlepage}
